\documentclass[11pt,a4paper]{article}
\usepackage[utf8]{inputenc}
\usepackage[T1]{fontenc}
\usepackage{lmodern}
\usepackage[french]{babel}
\usepackage{amsmath,amssymb,mathtools}
\usepackage{icomma}
\usepackage{xcolor}
\usepackage[most]{tcolorbox}
\usepackage{enumitem}
\usepackage[a4paper,margin=2.2cm,headheight=26pt,headsep=10pt]{geometry}
\usepackage{fancyhdr}
\usepackage[hidelinks]{hyperref}
\definecolor{primary}{RGB}{222,49,99}
\definecolor{primarydark}{RGB}{150,22,55}
\tcbset{boxbase/.style={enhanced,breakable,boxrule=0.9pt,arc=2.5pt,
  left=3mm,right=3mm,top=2.3mm,bottom=2.3mm,fonttitle=\bfseries,coltitle=white}}
\newtcolorbox[auto counter]{exo}[1][]{boxbase,colback=primary!4,colframe=primary,
  title={\textsc{Exercice}~\thetcbcounter\ifx\relax#1\relax\relax\else\textmd{ --- \itshape #1}\fi}}
\newcommand{\R}{\mathbb{R}}
\newcommand{\ds}{\displaystyle}
\pagestyle{fancy}\fancyhf{}
\fancyhead[L]{\small\textcolor{primarydark}{\textbf{Thème 3} — Factorisation et résolution}}
\fancyhead[R]{\small\textcolor{primarydark}{\textsc{Moyen}}}
\fancyfoot[C]{\small\thepage}
\renewcommand{\headrulewidth}{0.8pt}
\begin{document}

\begin{center}
{\LARGE\bfseries\color{primarydark} Niveau Moyen — Exercices}\\[2pt]
{\color{primary}\rule{0.45\textwidth}{1.2pt}}\\[3pt]
{\small\itshape Factorisations composées et méthode de Horner. Aucune calculatrice.}
\end{center}
\vspace{1mm}

\begin{exo}[mise en évidence avec puissances et parenthèses]
Factorise :
\begin{enumerate}[label=\alph*),leftmargin=2em,topsep=1pt,itemsep=1pt]
  \item $5\times 2^{2x}-3a\times 2^{2x+1}$ \qquad
  \item $e^{2x}+e^{x}$ \qquad
  \item $(x-1)(2x+3)-(x-1)^{2}$ \qquad
  \item $6x^{4}-4x^{3}+2x^{2}$
\end{enumerate}
\end{exo}

\begin{exo}[produits remarquables composés]
Factorise le plus complètement possible :
\begin{enumerate}[label=\alph*),leftmargin=2em,topsep=1pt,itemsep=1pt]
  \item $e^{2x}+2e^{x}+1$ \qquad
  \item $e^{4x}-1$ \qquad
  \item $5x^{2}-9$ \qquad
  \item $\bigl(\ln x^{2}\bigr)^{2}-3$ (avec $x>0$)
\end{enumerate}
\end{exo}

\begin{exo}[factoriser le trinôme, puis résoudre]
Résous dans $\R$ en factorisant d'abord :
\begin{enumerate}[label=\alph*),leftmargin=2em,topsep=1pt,itemsep=1pt]
  \item $2x^{2}-3x+1=0$ \qquad
  \item $3x^{2}-2\sqrt3\,x+1=0$ \qquad
  \item $x^{2}-x-6=0$
\end{enumerate}
\end{exo}

\begin{exo}[factorisation par la méthode de Horner]
Pour chaque polynôme, vérifie que $x=1$ est racine, puis factorise complètement :
\begin{enumerate}[label=\alph*),leftmargin=2em,topsep=1pt,itemsep=1pt]
  \item $P(x)=x^{3}-2x^{2}-x+2$ \qquad
  \item $P(x)=x^{3}-x^{2}+2x-2$
\end{enumerate}
\end{exo}

\begin{exo}[simplifier une fraction rationnelle]
Donne d'abord la condition d'existence (C.E.), puis simplifie :
\begin{enumerate}[label=\alph*),leftmargin=2em,topsep=1pt,itemsep=1pt]
  \item $\ds\frac{x^{2}-1}{x^{2}-3x+2}$ \qquad
  \item $\ds\frac{x^{2}-4}{x+2}$
\end{enumerate}
\end{exo}

\end{document}
